\section{独立監査 B：Class 6 量子作用素積からの時間 Lax 成分}
\subsection{問題設定と今回の計算の範囲}
Class 6は、最近接相互作用を持つ可積分なスピン二分の一の量子鎖の一つである。本報告の文献入力は、de Leeuw--Fontanella--Nieto Garcíaの論文\cite{DFN}の付録Aにある量子R行列と、比較用の連続運動方程式である。Class 6がeleven-vertex模型に対応することは既知であり\cite{CdL}、対応するLandau--Lifshitz型場の理論も既に研究されている\cite{LOZ}。従って「Class 6にLax表示がある」こと自体を新規な主張にはしない。

調べるのは、\emph{有限格子の量子Lax方程式の各作用素積をコヒーレント状態で評価したとき、通常の古典零曲率方程式にどう整理されるか}である。量子鎖から古典連続模型を得る方法論\cite{ADS}や、コヒーレント状態のsymbolの積を扱う方法\cite{KS}も既知である。本報告では、Class 6で必要な項を省かずに計算し、古典時間成分を先に知っていなくても局所補正を構成できるかを検査する。

計算の順序は、量子R行列、量子Hamiltonian、有限格子の時間Lax作用素、作用素積の期待値、連続極限、時間Lax成分の局所補正、独立な運動方程式との照合、の順とする。最後に得た古典時間成分を既知の式から差し引いて補正を\emph{定義}する方法は採らない。

本報告の式はすべてこの文書内で定義する。Class 5との比較は第\ref{audit6q:sec:comparison}節に限り、その比較に必要な規約もそこで再掲する。

\subsection{量子鎖、R行列、連続極限}
\subsubsection{空間と作用素の規約}
各格子点の量子空間を$\mathcal H_n=\mathbb C^2$、Lax作用素が付加的に作用する補助空間を$\mathcal V_0=\mathbb C^2$とする。添字$0$は補助空間、整数$n$は物理的な格子点を表す。二つの空間に作用する行列には、その空間の添字を付ける。例えば$R_{0n}$は$\mathcal V_0\otimes\mathcal H_n$に作用する。

$d$次単位行列を$\one d$、虚数単位を$\ii$、交換子を$[A,B]=AB-BA$と書く。二つの二次元空間を交換する置換作用素$P$と、Class 6の変形を表す行列$K$を、基底$(\uparrow\uparrow,\uparrow\downarrow,\downarrow\uparrow,\downarrow\downarrow)$で
\begin{equation}
P=\begin{pmatrix}1&0&0&0\\0&0&1&0\\0&1&0&0\\0&0&0&1\end{pmatrix},\qquad
K=\begin{pmatrix}0&1&1&0\\0&0&0&-1\\0&0&0&-1\\0&0&0&0\end{pmatrix}
\label{audit6q:eq:PK}
\end{equation}
と定義する。これらは
\begin{equation}
P^2=\one4,\qquad PK=KP=K,\qquad K^3=0
\label{audit6q:eq:PKalgebra}
\end{equation}
を満たす。ここでの$K$は量子空間上の四次行列であり、後に導入する古典スピンの異方性行列とは別物である。

\subsubsection{量子R行列とHamiltonian}
量子スペクトルパラメータを$u$、量子鎖の変形パラメータを$a$とする。文献\cite{DFN}の式(A.1)で交換相互作用の係数$a_1=1$としたR行列は
\begin{align}
R(u;a)&=2u\one4+P+a u(1+2u)K
 +\frac{a^2u^2(1+2u)^2}{2}K^2
\label{audit6q:eq:R}\\
&=\begin{pmatrix}
1+2u&au(1+2u)&au(1+2u)&-a^2u^2(1+2u)^2\\
0&2u&1&-au(1+2u)\\
0&1&2u&-au(1+2u)\\
0&0&0&1+2u
\end{pmatrix}.
\end{align}
以下、微分$\partial_u$では$a$を固定する。最近接の局所Hamiltonian$h$は
\begin{equation}
h=P\,\partial_uR(0;a)=2P+aK,
\qquad H_{\rm lat}=\sum_n h_{n,n+1}
\label{audit6q:eq:hlat}
\end{equation}
である。有限周期鎖を取る場合は添字を周期的に解釈する。時間を$t_{\rm lat}$とし、$\hbar=1$、Heisenberg方程式$\partial_{t_{\rm lat}}O=\ii[H_{\rm lat},O]$を用いる。

式\eqref{audit6q:eq:R}から直接、
\begin{equation}
R(0;a)=P,\qquad
R(u;a)R(-u;a)=(1-4u^2)\one4
\label{audit6q:eq:inverse}
\end{equation}
が得られる。後者は逆行列を与える代数的関係であり、Hermitian共役を含む物理的なユニタリティ条件ではない。逆行列を使う有限格子計算では$u\ne\pm\tfrac12$とする。

\subsubsection{格子間隔に関する展開}
格子間隔を$\epsilon$、位置を$x_n=n\epsilon$とする。連続理論の変形パラメータ$\alpha$、連続スペクトルパラメータ$\lambda\ne0$、連続時間$T$を、
\begin{equation}
a=\epsilon^2\alpha,\qquad u=\frac{\lambda}{\epsilon},\qquad
T=\epsilon^2t_{\rm lat}
\label{audit6q:eq:scaling}
\end{equation}
によって定義する。長さの単位を戻すなら、$u,a$を無次元として$\lambda$は長さ、$\alpha$は長さの逆二乗、$T$は長さの二乗の次元を持つ。

量子空間Lax作用素を$L_n=R_{0n}(u;a)$とする。連続極限で単位行列に近づく規格化を選び、
\begin{equation}
\widehat L_n=\frac{1}{2u}L_n
\label{audit6q:eq:Lnorm}
\end{equation}
と定義する。格子間隔の一次、二次、三次の作用素係数を、それぞれ$\ell_1,\ell_2,\ell_3$と書く。式\eqref{audit6q:eq:R}の代入だけで、厳密な有限多項式
\begin{align}
\widehat L&=\one4+\epsilon\ell_1+\epsilon^2\ell_2+\epsilon^3\ell_3,
\label{audit6q:eq:Lseries}\\
\ell_1&=\frac{P}{2\lambda}+\alpha\lambda K+\alpha^2\lambda^3K^2,
\label{audit6q:eq:l1}\\
\ell_2&=\frac{\alpha}{2}K+\alpha^2\lambda^2K^2,
\qquad
\ell_3=\frac{\alpha^2\lambda}{4}K^2
\label{audit6q:eq:l23}
\end{align}
を得る。Class 6では二次係数$\ell_2$を落とせない。特に
\begin{equation}
\ell_1^2=\frac{1}{4\lambda^2}\one4+2\ell_2
\label{audit6q:eq:l1square}
\end{equation}
が成り立ち、一次係数の作用素としての二乗はスカラー行列ではない。

有限格子の時間成分で必要なスペクトル微分は、規格化した作用素を$u$で微分してから式\eqref{audit6q:eq:scaling}を入れると
\begin{equation}
\left.\partial_u\widehat L\right|_{a\ {\rm fixed}}
=\epsilon^2\partial_\lambda\ell_1
 +\epsilon^3\partial_\lambda\ell_2
 +\epsilon^4\partial_\lambda\ell_3
\label{audit6q:eq:spectralderivative}
\end{equation}
となる。右辺の$\lambda$微分では$\alpha$を固定する。

\subsection{コヒーレント状態と空間Lax行列}
\subsubsection{期待値で何を残すか}
古典スピン場を$\Svec(T,x)=(S^1,S^2,S^3)^{\mathsf T}$とし、単位長の拘束$\Svec^{\mathsf T}\Svec=1$を課す。内積は複素共役を取らない双線形形式である。Pauli行列は
\begin{equation}
\sigma_1=\begin{pmatrix}0&1\\1&0\end{pmatrix},\quad
\sigma_2=\begin{pmatrix}0&-\ii\\\ii&0\end{pmatrix},\quad
\sigma_3=\begin{pmatrix}1&0\\0&-1\end{pmatrix}
\label{audit6q:eq:Pauli}
\end{equation}
とする。スピン・コヒーレント状態の密度行列を
\begin{equation}
\rho(\Svec)=\frac12\left(\one2+\sum_{a=1}^3S^a\sigma_a\right)
\label{audit6q:eq:rho}
\end{equation}
と書く。実単位球面上では通常の純粋状態である。非Hermitianな模型の式を複素化した場へ拡張した場合には、$\rho^2=\rho$と$\tr\rho=1$は保たれるが、正値な密度行列であるとは主張しない。

作用素$B$の上付き矢印$B^\downarrow$は、\emph{物理的な量子空間だけで期待値を取り、補助空間の行列を残す}操作、すなわちlower symbolを意味する。一格子点では
\begin{equation}
B^\downarrow(\Svec)=\tr_n[(\one2\otimes\rho(\Svec))B_{0n}],
\label{audit6q:eq:lower}
\end{equation}
二格子点では
\begin{equation}
B^\downarrow(\Svec_1,\Svec_2)
=\tr_{1,2}[(\one2\otimes\rho(\Svec_1)\otimes\rho(\Svec_2))B]
\label{audit6q:eq:lower2}
\end{equation}
である。$\tr_{1,2}$は物理空間1と2の部分トレースを表す。

異なる物理空間の状態が直積でも、同じ格子点を共有する作用素積について、期待値を先に取った行列の積へ置き換えてはならない。例えば一格子点の作用素を
$A=A_0\otimes\one2+\sum_a A_a\otimes\sigma_a$、
$B=B_0\otimes\one2+\sum_b B_b\otimes\sigma_b$と展開する。ここで$A_a,B_b$は補助空間上の行列である。Kronecker記号を$\delta_{ab}$、$\varepsilon_{123}=1$の完全反対称記号を$\varepsilon_{abc}$とすると、
\begin{equation}
(AB)^\downarrow-A^\downarrow B^\downarrow
=\sum_{a,b=1}^3 A_a B_b
\left(\delta_{ab}-S^aS^b+\ii\sum_{c=1}^3\varepsilon_{abc}S^c\right).
\label{audit6q:eq:star}
\end{equation}
この式はスピン二分の一に対して厳密であり、補助行列の積の順序も保っている。反対称記号$\varepsilon_{abc}$と格子間隔$\epsilon$は別の記号である。

\subsubsection{空間Lax行列の明示式}
複素スピン成分を$S^+=S^1+\ii S^2$、$S^-=S^1-\ii S^2$と定義する。これらは一般の複素化した場では互いの複素共役とは限らず、拘束は$S^+S^-+(S^3)^2=1$となる。

連続空間Lax行列$U(T,x;\lambda)$を、規格化した作用素の一次係数のlower symbolとして
\begin{equation}
U=\ell_1^\downarrow
=\frac1{4\lambda}
\begin{pmatrix}
1+S^3+2\alpha\lambda^2S^+&
S^-+4\alpha\lambda^2S^3-4\alpha^2\lambda^4S^+\\
S^+&1-S^3-2\alpha\lambda^2S^+
\end{pmatrix}
\label{audit6q:eq:U}
\end{equation}
と定める。従って$\widehat L_n^\downarrow=\one2+\epsilon U+O(\epsilon^2)$である。トレースを除いた空間行列を$U_0$と書き、
\begin{equation}
U_0=U-\frac1{4\lambda}\one2
\label{audit6q:eq:U0}
\end{equation}
と定義する。$U$と$U_0$は以下で区別する。

後の式を短くするため、三成分ベクトルをトレース零の二次行列へ写す線形写像$\Qmap$を導入する。任意の$\vvec=(v^1,v^2,v^3)^{\mathsf T}$に対し、$v^\pm=v^1\pm\ii v^2$として
\begin{equation}
\Qmap(\vvec)=\frac1{4\lambda}
\begin{pmatrix}
v^3+2\alpha\lambda^2v^+&v^-+4\alpha\lambda^2v^3-4\alpha^2\lambda^4v^+\\
v^+&-v^3-2\alpha\lambda^2v^+
\end{pmatrix}.
\label{audit6q:eq:Qmap}
\end{equation}
この写像は保存量ではない。特に$U_0=\Qmap(\Svec)$、$\partial_TU=\Qmap(\partial_T\Svec)$である。

\subsection{有限格子の時間成分と必要な展開次数}
\subsubsection{Sutherland relationからの構成}
有限格子の時間Lax作用素は、隣接する二つの物理格子点と補助空間に作用する。まず$L_n=R_{0n}$に対して、式\eqref{audit6q:eq:R}を代入するとSutherland relation
\begin{equation}
[h_{n-1,n},L_nL_{n-1}]
=L_n\partial_uL_{n-1}-(\partial_uL_n)L_{n-1}
\label{audit6q:eq:Sutherland}
\end{equation}
が厳密に成立する。これに基づいて時間作用素$\mathcal A_n$を
\begin{equation}
\mathcal A_n=\ii h_{n-1,n}
-\ii L_n^{-1}h_{n-1,n}L_n-\ii L_n^{-1}\partial_uL_n
\label{audit6q:eq:Atime}
\end{equation}
と取ると、
\begin{equation}
\partial_{t_{\rm lat}}L_n
=\ii[H_{\rm lat},L_n]
=\mathcal A_{n+1}L_n-L_n\mathcal A_n
\label{audit6q:eq:discreteZC}
\end{equation}
を満たす。ここでは有限格子の作用素等式そのものを検査する。

規格化$\widehat L_n=L_n/(2u)$を使うと、共通のスカラー時間成分を除くために
\begin{align}
\widehat{\mathcal A}_n&=\mathcal A_n+\frac{\ii}{u}\one8
\nonumber\\
&=-\ii\widehat L_n^{-1}
\left([h_{n-1,n},\widehat L_n]+\partial_u\widehat L_n\right)
\label{audit6q:eq:Anorm}
\end{align}
と定義できる。添字に対応する恒等作用は省略している。スカラーは両隣で相殺されるため、式\eqref{audit6q:eq:discreteZC}はハット付きの作用素でも成立する。

\subsubsection{格子間隔二次までの時間作用素}
この小節ではテンソル積の順序を$(0,1,2)=($補助空間、左格子点、右格子点$)$に固定する。変形のない局所Hamiltonianを$h_0=2P_{12}$、時間作用素の一次・二次係数を$A_1,A_2$として
\begin{equation}
\widehat{\mathcal A}=\epsilon A_1+\epsilon^2 A_2+O(\epsilon^3)
\label{audit6q:eq:Aseries}
\end{equation}
と展開する。式\eqref{audit6q:eq:Anorm}の直接展開から
\begin{align}
A_1&=-\ii[h_0,\ell_{1,02}],
\label{audit6q:eq:A1}\\
A_2&=\ii\ell_{1,02}[h_0,\ell_{1,02}]
-\ii[h_0,\ell_{2,02}]-\ii\partial_\lambda\ell_{1,02}
\label{audit6q:eq:A2}
\end{align}
となる。$h=h_0+\epsilon^2\alpha K_{12}$なので、Hamiltonianの変形はこの展開では三次以降の時間作用素に入る。その寄与を物理的に無視しているわけではない。

